\documentclass[11pt]{article}
\usepackage[utf8]{inputenc}
\usepackage{cmap}
\usepackage[T2A]{fontenc}
\usepackage[russian]{babel}
\usepackage{amsmath}
\usepackage{graphicx}
\usepackage{geometry}
\geometry{a4paper, margin=2.5cm}

\date{}

\begin{document}

\begin{titlepage}

\begin{center}
\large
\textbf{ОПРЕДЕЛЕНИЕ НИЖНЕЙ ГРАНИЦЫ ОБЛАЧНОСТИ\\ ПО СТЕРЕОИЗОБРАЖЕНИЯМ С ИСПОЛЬЗОВАНИЕМ\\ ГЛУБОКОГО СОПОСТАВЛЕНИЯ ПРИЗНАКОВ\\ И ВРЕМЕННОЙ ФИЛЬТРАЦИИ}

\vspace{10mm}

\large
\textbf{Устимов И.А.$^{(1)}$, Постыляков О.В.$^{(2)}$, Чуличков А.И.$^{(1)}$, Криницкий М.А.$^{(3,4)}$, Борисов М.А.$^{(4)}$, Белкова Е.А.$^{(4)}$, Лулукин В.О.$^{(2)}$}

\vspace{5mm}

\normalsize
$^{(1)}$~Кафедра математического моделирования и информатики, Физический факультет, МГУ имени М.В.~Ломоносова, Москва, 119991, Россия\\
$^{(2)}$~Институт физики атмосферы имени А.М.~Обухова РАН, Москва, 119017, Россия\\
$^{(3)}$~Московский физико-технический институт (НИУ), Долгопрудный, 141701, Россия\\
$^{(4)}$~Институт океанологии имени П.П.~Ширшова РАН, Москва, 117997, Россия\\
\texttt{ustimov.ia@physics.msu.ru}

\vspace{20mm}

\textbf{Аннотация}
\end{center}

\vspace{5mm}

Представлен пассивный стереоскопический метод оценки высоты нижней границы облачности (НГО) по снимкам с наземных fisheye-камер. Основу метода составляет LoFTR (Local Feature TRansformer)~--- детекторно-независимая архитектура глубокого обучения на основе Vision Transformer, которая строит плотные субликсельные соответствия непосредственно по паре изображений без отдельного этапа выделения ключевых точек. Диспарантности кластеризуются с помощью алгоритма HDBSCAN для разделения облачных слоёв и отбраковки выбросов. Трёхстадийный пространственный фильтр устраняет физически неправдоподобные кластеры, а адаптивный фильтр Калмана с зависящей от невязки ковариацией шума измерений подавляет межкадровые скачки высоты до 500~м за 10~с, сохраняя реальные вариации высоты.

Валидация выполнена по данным совместно расположенного лазерного высотомера ДОЛ-2 за 5-часовой дневной цикл (23 мая 2025~г.). После временной фильтрации и разбиения на 30-минутные окна с индивидуальной динамической временной синхронизацией средняя корреляция Пирсона с высотомером составляет $r = 0{,}67$; в 72\% окон (13 из 18) $r > 0{,}5$, максимальное значение~--- $r = 0{,}94$. Фильтр Калмана снижает средний межкадровый перепад высоты в 4{,}3 раза и повышает 5-часовую автокорреляцию ряда с $0{,}11$ до $0{,}90$. Систематическое смещение между оценками стереосистемы и высотомера объясняется принципиальным различием измеряемых величин: геометрической границы облака (по рассеянному солнечному свету) и оптической границы (по обратному рассеянию лазерного импульса).

\vspace{10mm}

\noindent\textbf{Ключевые слова:} высота нижней границы облачности, стереозрение, глубокое сопоставление признаков, LoFTR, фильтр Калмана, HDBSCAN, пространственная фильтрация.

\vspace{10mm}

\begin{flushright}
\normalsize
Устимов Иван Андреевич\\
\textit{ustimov.ia@physics.msu.ru}
\end{flushright}

\end{titlepage}

\end{document}
